

\section{Finite Semigroups and the Road Coloring Problem}  
\label{sec:4.2}

A finite semigroup is a finite set, $S$, which is closed under a binary associative operation. 
For any subset $T$ of $S$, we will write $E(T)$ to refer to the set of idempotents in $T$. 
Let $G=(V,\mathcal{E})$ be a strongly connected digraph with $b$ vertices having adjacency matrix $A$. 
We will assume that the outdegree of each vertex of $G$ or, equivalently, the sum of each  row of $A$ is some constant $d$.

\begin{definition}	\label{organized:coloring-semigroup-definition}
A coloring of $G$ is a decomposition of $A$ into $d$ binary (that is, $0-1$)
stochastic matrices $R_1,R_2,\dots,R_d$ such
that $A=R_1+\cdots+R_d$.

\vphantom{...}

Intuitively, one can think of a matrix $R_i$ as assigning the $i$th of $d$
colors to the edge $(j,k)$ if $(R_i)_{j,k}=1$. Each 
$b\times b$ binary stochastic matrix $R$ can also be thought of as a
function on the vertices $\{1,2,\dots,n\}$ where
$R_{j,k}=1$ iff $jR=k$. As a notational convention, writing functions on the
right allows us to be indifferent to whether we
consider $jR=k$ as a function mapping $j$ to $k$ or $R$ as a binary
stochastic matrix with $j$ a row vector having a one in 
the $j$th position and zeros elsewhere. 

Let $S=\left\langle R_1,R_2,\dots, R_d\right\rangle$ be the semigroup
generated by $R_1,R_2,\dots,R_d$ under 
matrix multiplication (or composition). Clearly, $S$ is a finite semigroup.
We will refer to $S$ as a coloring semigroup. The
following theorem is extremely important in all that follows. 
\end{definition}

\begin{theorem}[\cite{CP}]	\label{thm:4.1} 
Let $S$ be a finite semigroup. Then $S$ contains a minimal ideal $\mathcal{K}$
called the kernel which is a disjoint union of
isomorphic group. In fact, $\mathcal{K}$ is isomorphic to $\mathcal{X} \times \mathcal{G}  \times \mathcal{Y}$ where if $e\in E(S)$, then
$e\mathcal{K}e$ is a group and
\[
\mathcal{X} = E( \mathcal{K}e)\qquad  \mathcal{G} = e\mathcal{K}e \qquad  \mathcal{Y} = E(e\mathcal{K}
)
\]
and if $\left( x_1,g_1,y_1\right) $, $\left( x_2,g_2,y_2\right)  \in \mathcal{X} \times \mathcal{G} \times \mathcal{Y}$
then 
\[
(x_1, g_1, y_1) (x_2, g_2, y_2) = (x_1, g_1( y_1x_2) g_2, y_2)
\]
The product structure $\mathcal{X} \times \mathcal{G} \times \mathcal{Y}$ is sometimes called a Rees product and any semigroup that has 
a Rees product is called completely simple. 
Thus the kernel of a finite semigroup is always completely simple.
\end{theorem}

The next theorem is extremely useful for characterizing the elements of the
kernel. It is an elementary implication of a more 
general result by Clark \cite{Cl}.

\begin{theorem}		\label{thm:4.2}  
Let $S$ be a finite semigroup of matrices. 
Then the kernel of $S$ is the set of matrices with minimal rank.
\end{theorem}

\begin{example}		\label{exm:4.1}
Let us reconsider the road coloring problem using the following example.
Suppose the adjacency matrix of a graph is given by 
\[
  A=\begin{bmatrix}
	0 &2 & 0 \\
	1 & 0 & 1 \\
	1 & 0 & 1
  \end{bmatrix}.
\]
Let
\[
\begin{alignedat}{2}
  R_1 &=\bmatrix 0 & 1 & 0 \\ 1 & 0 & 0 \\ 1 & 0 & 0\endbmatrix & \qquad 
  B_1 &=\bmatrix 0 & 1 & 0 \\ 0 & 0 & 1 \\ 0 & 0 & 1\endbmatrix \\
\\
  R_2 &=\bmatrix 0 & 1 & 0 \\ 1 & 0 & 0 \\ 0 & 0 & 1\endbmatrix & \qquad
  B_2 &=\bmatrix 0 & 1 & 0 \\ 0 & 0 & 1 \\ 1 & 0 & 0\endbmatrix
\end{alignedat}
\]
Then $A=R_1+B_1=R_2+B_2$.

Let $S_1=\left\langle R_1,B_1\right\rangle$ and $S_2=\left\langle
R_2,B_2\right\rangle$. 
Notice that both $R_2$ and $B_2$ are permutation matrices, thus $S_2$ is a group and contains only matrices of rank three. 
However, 
\[
  \left( B_1\right) ^2=\begin{bmatrix}
	0 & 0 & 1 \\
	0 & 0 & 1 \\
	0 & 0 & 1
  \end{bmatrix}.
\]

Therefore, the kernel of $S_1$ is a set of rank one matrices. In particular,
$\left( B_1\right) ^2$ is a synchronizing instruction. 
Following the blue road twice will take each person to building three. Thus
we can make the following observations from 
this example:

\vphantom{...}
%
\begin{enumerate}
\item[(1)] Rank one matrices (equivalently, constant functions) have a
natural interpretation as synchronizing instructions

\item[(2)] Each graph will, in general, have many colorings and not
every coloring will produce a synchronizing instruction.
\end{enumerate}
%
\end{example}

The existence of a synchronizing instruction has ramifications for the
``period'' of a given graph. We first define the period 
of a vertex. 

\begin{definition}		\label{defn:4.2}
For each $v \in V$, define the period of $v$ as 
\[
  \per(v)=\gcd\{|W|:W\text{ is a walk from }v\text{ to }v\}.
\]

If the $\per(v) = 1$, then $v$ is called aperiodic. 
It is a well known fact in Markov chain theory that the period of a vertex 
is a so-called class property. 
That is, each vertex in a given strongly connected component (irreducible component, in Markov chain terminology) has the same period. 
Our assumption of strong connectedness for the graph allows us to speak 
of the period of the graph. 
The following theorem was originally proven in
\cite{AGW}.
\end{definition}

\begin{theorem}		\label{thm:4.3}
If $G$ has a synchronizing instruction, then $G$ is aperiodic.
\end{theorem}

\begin{proof}
Let $S=\left\langle R_1,R_2,\dots,R_d\right\rangle$ be a coloring semigroup
for $G$ that contains a synchronizing instruction $W=M_1M_2\dots M_k$ where each $M_i$ is one of the colors $R_j$. 
Then $W$ is a rank one binary stochastic matrix with one column,  say $w$, which is nonzero. 
Then for each vertex $v$ in $V$, $vW=w$. 
In particular, $W$ represents a walk of length $k$ from $w$ to itself. Suppose $u=wR_1$.
Now  $wR_1W= uW = w$, thus $R_1W$ is a walk of length $k+1$ from $w$ to itself. Since the period of $w$ is $\gcd(k,k+1)$, $G$ is aperiodic.
\qed
\end{proof}

Thus aperiodicity is a necessary condition for the existence of a synchronizing instruction. 
The road coloring problem seeks to determine if aperiodicity is also a sufficient condition.

\runinhead{Road Coloring Conjecture}
Let $G=(V,\mathcal{E})$ be a $d$-out, strongly connected digraph with adjacency matrix $A$. 
Let $\left\{\left\{R_{i_1},R_{i_2},\dots,R_{i_d}\right\}:A=R_{i_1}+\cdots
+R_{i_d}\right\}$ be the set of all colorings of $G$ and let $S_i=\left\langle R_{i_1},R_{i_2},\dots,R_{i_d}\right\rangle$ be the associated coloring semigroups. 
If $G$ is aperiodic, then for some $i$, $S_i$ contains a synchronizing instruction.
